\subsubsection{有限维 Fredholm 行列式的相位梳齿刚性：极点晶格与计数上界}\label{subsubsec:operator-finite-dimensional-phase-comb}
注 \ref{rem:finite-pole-lattice-barrier} 已给出有限维实现的结构性门槛：当 \(\widehat\zeta\) 由有限矩阵的行列式给出时，其极点在 \(s\)-平面上只能呈现为有限条竖直直线上的等差梳齿。
本段将该门槛表述为可直接引用的计数不等式，并给出临界线达界谱点对应的同余类上界。

\begin{theorem}[有限维 \texorpdfstring{$\widehat\zeta$}{zeta-hat} 的极点梳齿结构与高度窗计数上界]\label{thm:operator-finite-dimensional-zhat-pole-counting}
令 \(M\in\CC^{d\times d}\) 为矩阵，设 \(\lambda:=\rho(M)>0\)。
定义
\[
\widehat\zeta_M(s):=\det\!\bigl(I-\lambda^{-s}M\bigr)^{-1}.
\]
则对每个非零特征值 \(\mu\in\mathrm{spec}(M)\setminus\{0\}\)，\(\widehat\zeta_M\) 的极点由方程 \(\lambda^{-s}\mu=1\) 给出，亦即存在参数化
\[
\boxed{\;
s
=\frac{\log|\mu|}{\log\lambda}
\;+\;
\mathrm{i}\,\frac{\arg\mu+2\pi k}{\log\lambda},
\qquad k\in\ZZ.\;}
\]
因此极点集合分解为有限条竖直直线上的等差点列，其虚部公差为
\[
\Delta:=\frac{2\pi}{\log\lambda}.
\]
进一步，任取 \(H>0\) 与任意固定竖直线 \(\Re(s)=\sigma\)。
令 \(N_{\sigma,H}(\widehat\zeta_M)\) 表示该直线内高度窗 \(|\Im(s)|\le H\) 中 \(\widehat\zeta_M\) 的极点数（按重数计）。
则有上界
\[
\boxed{\;
N_{\sigma,H}(\widehat\zeta_M)\ \le\ d\Bigl(1+\frac{H\log\lambda}{\pi}\Bigr).\;}
\]
特别地，若仅计入临界线 \(\Re(s)=\tfrac12\) 上的极点，则
\[
\boxed{\;
N_{\tfrac12,H}(\widehat\zeta_M)\ \le\ d_{1/2}\Bigl(1+\frac{H\log\lambda}{\pi}\Bigr),\qquad
d_{1/2}:=\#\{\mu:\ |\mu|=\lambda^{1/2}\}\le d.\;}
\]
\end{theorem}
\begin{proof}
极点条件为 \(\det(I-\lambda^{-s}M)=0\)，等价于存在 \(\mu\in\mathrm{spec}(M)\setminus\{0\}\) 使 \(\lambda^{-s}\mu=1\)。
取对数得到
\[
-s\log\lambda+\log\mu=2\pi\mathrm{i}k,\qquad k\in\ZZ,
\]
写 \(\log\mu=\log|\mu|+\mathrm{i}\arg\mu\) 即得所示参数化与公差 \(\Delta\)。
\par
对固定的 \(\mu\)，不等式 \(|\Im(s)|\le H\) 等价于
\[
\left|\frac{\arg\mu+2\pi k}{\log\lambda}\right|\le H,
\]
其满足的整数 \(k\) 个数至多为 \(1+H\log\lambda/\pi\)。
对所有特征值按代数重数求和，即得 \(N_{\sigma,H}\) 的上界；临界线版本仅对满足 \(\log|\mu|/\log\lambda=\tfrac12\) 的谱点求和。证毕。
\end{proof}

\begin{corollary}[虚部同余类数目上界：有限维模型的相位梳齿刚性]\label{cor:operator-phase-comb-congruence-classes}
沿用定理 \ref{thm:operator-finite-dimensional-zhat-pole-counting} 的记号，令 \(\Delta=2\pi/\log\lambda\)。
则对任意固定竖直线 \(\Re(s)=\sigma\)，该直线上极点虚部集合在模 \(\Delta\) 的意义下至多落在 \(d\) 个同余类（更精确地：至多落在有限集合 \(\{\arg\mu/\log\lambda:\ \mu\neq 0,\ \log|\mu|/\log\lambda=\sigma\}\) 所给出的平移类上）。
因此，若某一目标谱几何要求在同一竖直线内出现超过 \(d\) 个互异的 \((\bmod\,\Delta)\) 虚部类，则任何维数 \(\le d\) 的有限维核都不可能实现之；要获得更丰富的临界线几何/统计，必须引入状态数增长的极限过程或进入无穷维转移算子的 Fredholm 行列式口径（见注 \ref{rem:finite-pole-lattice-barrier}）。
\end{corollary}
\begin{proof}
由定理 \ref{thm:operator-finite-dimensional-zhat-pole-counting} 的参数化，虚部必形如 \((\arg\mu+2\pi k)/\log\lambda\)。
模 \(\Delta\) 仅保留 \(\arg\mu/\log\lambda\) 的有限取值；而不同竖直线对应不同的 \(|\mu|\)。
对固定 \(\sigma\) 可出现的 \(\mu\) 数目至多为 \(d\)（按代数重数计），故同余类数目至多为 \(d\)。证毕。
\end{proof}

\endinput
